\documentclass[10pt]{article}
\usepackage[margin=0.75in]{geometry}
\usepackage{amsmath,amssymb,amsthm,microtype}
\newtheorem{theorem}{Theorem}
\newcommand{\R}{\mathcal R}
\newcommand{\A}{\mathcal A}
\begin{document}
\begin{center}
{\Large Combining Two Smaller Transfer Instances After an Inverse Step}\par\medskip
Collatz proof project\quad---\quad September 5, 2026
\end{center}
Let $T$ be the shortcut Collatz map and
$\A(u)=[\R(3u+2)\Rightarrow\R(27u+20)]$, where $\R$ means reaching one.
The three-step inverse template at parameter 15 supplies an auxiliary orbit
but fails its single-merge criterion. A second smaller transfer instance
completes the certificate on an entire progression. The recursive premises
remain explicit, and global coverage remains open.

\begin{theorem}[Two-premise induction step]
For any natural $Q$, set
\[
 u=15+9\cdot2^{40}Q,\qquad
 v=13+8\cdot2^{40}Q,\qquad
 w=2+3^{16}2^{15}Q.
\]
Then $v<u$, $w<u$, and
\begin{align*}
 T^3(3v+2)&=3u+2,\\
 T^{28}(27v+20)&=3w+2,\\
 T^{15}(27w+20)&=T^{40}(27u+20)=2+3^{25}Q.
\end{align*}
Consequently $\A(v)\land\A(w)\Rightarrow\A(u)$.
\end{theorem}
\begin{proof}
The four base paths have endpoints and odd-step counts
\[
\begin{array}{c|c|c}
\text{path}&\text{endpoint}&\text{odd count}\\\hline
T^3(41)&47&2\\
T^{28}(371)&8&14\\
T^{15}(74)&2&6\\
T^{40}(425)&2&20
\end{array}
\]
Their input increments are respectively
$2^3(3\cdot2^{40}Q)$,
$2^{28}(27\cdot2^{15}Q)$,
$2^{15}(3^{19}Q)$, and $2^{40}(243Q)$.
Binary-cylinder transport gives every displayed identity.
For strict decrease, $u-v=2+2^{40}Q>0$, and the coefficient comparison
for $w$ reduces to $3^{14}<2^{25}$, alongside $2<15$.

Assume $\R(3u+2)$. The first identity transports convergence backward to
$3v+2$. Apply $\A(v)$ and advance 28 steps to $3w+2$. Apply $\A(w)$
and advance 15 steps to the target's 40-step endpoint. Transport convergence
back to $27u+20$. Lean checks the geometry and this argument for arbitrary $Q$.
\end{proof}

\paragraph{How the second bridge was found.}
For $u=9a+6$, the inverse template uses the smaller parameter $8a+5$.
After its bridge and three forward steps, the remaining pair is
$243a+175$ and $243a+182$ with the same slope.
At $a=1$, the auxiliary base value reaches 8 after another 25 steps.
This exposes the second recursive base parameter $(8-2)/3=2$.
Its larger partner is 74, whose 15-step endpoint supplies the required merge.
Thus the initial single-template failure is not a failure of all configurations
of smaller transfer instances.

\paragraph{Finite search and verification.}
The reproducible search allows one extra forward bridge after the inverse
template. On $a=0,\ldots,255$, through comparison depth 60, it finds
142 conditional certificates and leaves 114 at the depth limit; there are
no state-cap outcomes. These are finite sample counts, not residue-density
or global-coverage claims. Three tests replay every successful sample at
$Q=0,1,7$, check that the parameter-15 certificate needs the extra bridge,
and distinguish the explicit state cap. The saved report contains the
counts and the parameter-15 example.

\paragraph{Scope and remaining goal.}
The two public declarations in \texttt{Collatz.TwoEarlyBridges} prove the
uniform strict decreases, all orbit identities, and the conditional implication.
The search census is not kernel certified. No convergence premise for the
unknown lifted target is assumed. Nevertheless this one progression does
not treat every Mersenne parameter $2^k-1$; it contains the base case 15,
not a uniform family of those exponents. The missing global argument still
requires coverage of all parameters by decreasing rules or another method.
No mathematical priority claim or proof of Collatz is asserted.
\end{document}
